% -----------------------------------------------------------------------------
% -----------------------------------------------------------------------------

\begin{exbox}[QPINEM 真空损失]
量子交换传播子~\eqnrefs{eq:qpinem_S} 对 Fock 输入的弱交换条件为 $|g_{\mathrm Q}|^2(n+1)\ll1$。在实际占据范围内满足该条件后，展开为
\begin{equation*}
 S_Q
 =1+g_{\mathrm Q}\hat b^\dagger\hat a
 -g_{\mathrm Q}^*\hat b\hat a^\dagger+O(|g_{\mathrm Q}|^2).
\end{equation*}
对初态 $|0,n\rangle$，第一项保持电子能量不变；第二项湮灭一个光子并把电子推到 $\ell=+1$；第三项产生一个光子并把电子推到 $\ell=-1$。因此
\begin{align}
 P_{+1|n}&=|g_{\mathrm Q}|^2n+O(|g_{\mathrm Q}|^4),\notag\\
 P_{-1|n}&=|g_{\mathrm Q}|^2(n+1)+O(|g_{\mathrm Q}|^4).
 \label{eq:qpinem-gain-loss-n-nplus1-dp65}
\end{align}
这里的 $+1$ 是电子\emph{增益}一个光子能量，$-1$ 是电子\emph{损失}一个光子能量。两者的差
\begin{equation}
 \boxed{P_{-1|n}-P_{+1|n}=|g_{\mathrm Q}|^2+O(|g_{\mathrm Q}|^4)}
 \label{eq:qpinem-spontaneous-difference-dp65}
\end{equation}
与 $n$ 无关，正是玻色对易关系 $[\hat a,\hat a^\dagger]=1$ 留下的真空项。

例如取 $|g_{\mathrm Q}|=0.10$：真空 $n=0$ 时
\begin{equation*}
 P_{+1|0}=0,
 \qquad
 P_{-1|0}\simeq0.010.
\end{equation*}
所以“没有入射光子”并不意味着电子不能损失能量；电子仍可向真空模式自发发射一个量子。对 $n=1$，增益和损失概率分别约为 $0.010$ 与 $0.020$；对 $n=10$，则约为 $0.10$ 与 $0.11$。只有同时维持 $|g_{\mathrm Q}|^2(n+1)\ll1$ 时，增大 $n$ 才使这两条弱交换率的相对差异按 $1/n$ 减小。这只检验增益--损失对称，不能据此把 Fock 输入的电子约化态视为经典相干调制后的纯态。


\medskip

\eqnrefs{eq:qpinem_S} 给出单模 QPINEM 的完整联合散射算符。在 $|g_{\mathrm Q}|\ll1$ 的弱交换极限，散射算符对初态 $|0,n\rangle$ 的一阶振幅只含一个光子升降算符，因此 Fock 梯矩阵元直接给出 $\sqrt n$ 与 $\sqrt{n+1}$ 的玻色增强。概率保留到 $O(|g_{\mathrm Q}|^2)$ 即得到\eqnrefs{eq:qpinem-gain-loss-n-nplus1-dp65}；电子双边平移仍要求实际占据通道远离谱下界。


由指数展开
\begin{equation*}
 S_Q=1+X+O(X^2),
 \qquad
 X=g_{\mathrm Q}\hat b^\dagger\hat a-g_{\mathrm Q}^*\hat b\hat a^\dagger.
\end{equation*}
作用在 $|0,n\rangle$ 上：
\begin{align*}
 \hat b^\dagger\hat a|0,n\rangle
 &=\sqrt n\,|+1,n-1\rangle,\\
 \hat b\hat a^\dagger|0,n\rangle
 &=\sqrt{n+1}\,|-1,n+1\rangle.
\end{align*}
因此
\begin{equation*}
 S_Q|0,n\rangle
 =|0,n\rangle
 +g_{\mathrm Q}\sqrt n\,|+1,n-1\rangle
 -g_{\mathrm Q}^*\sqrt{n+1}\,|-1,n+1\rangle
 +O(|g_{\mathrm Q}|^2).
\end{equation*}
两个末态的电子边带和光子数都互相正交，故到这一阶不存在交叉项。振幅取模平方即得
\begin{equation*}
 P_{+1|n}=|g_{\mathrm Q}|^2n,
 \qquad
 P_{-1|n}=|g_{\mathrm Q}|^2(n+1).
\end{equation*}
相减后只剩 $|g_{\mathrm Q}|^2$。从算符层面看，这个 $1$ 来自
\begin{equation*}
 \langle n|\hat a\hat a^\dagger|n\rangle
 =\langle n|(\hat a^\dagger\hat a+1)|n\rangle=n+1,
\end{equation*}
所以它不能由经典复振幅替代；经典场极限只是在 $n\gg1$ 时把相对真空修正压到 $1/n$。
\end{exbox}



\begin{exbox}[相同平均光子数的不同电子谱]
现在固定三种光场都具有同一平均光子数
\begin{equation*}
 \bar n=\langle\hat n\rangle=2,
\end{equation*}
但分别取 Fock 态 $|2\rangle$、平均数为 2 的相干态和平均数为 2 的单模热态。它们的一阶平均完全相同，因此在 $|g_{\mathrm Q}|=0.10$ 的弱耦合下都有
\begin{equation*}
 P_{+1}\simeq|g_{\mathrm Q}|^2\bar n=0.020.
\end{equation*}
但二阶增益边带由\eqnrefs{eq:qpinem-weak-factorial-moment-dp19} 控制：
\begin{equation*}
 P_{+2}
 \simeq\frac{|g_{\mathrm Q}|^4}{4}
 \langle\hat n(\hat n-1)\rangle.
\end{equation*}
三种光场的结果是
\begin{center}
\begin{tabular}{c|c|c|c}
 光场 & $\langle\hat n(\hat n-1)\rangle$ & $P_{+2}$ & $g^{(2)}(0)$\\\hline
 $|2\rangle$ & 2 & $5.0\times10^{-5}$ & $1/2$\\
 相干态，$|\alpha|^2=2$ & 4 & $1.0\times10^{-4}$ & $1$\\
 单模热态，$\bar n=2$ & 8 & $2.0\times10^{-4}$ & $2$
\end{tabular}
\end{center}
也就是说，三种场的平均光子数完全相同，电子的一阶增益概率也相同，但二阶边带相差四倍。用只包含电子计数的组合
\begin{equation}
 \boxed{
 \frac{4P_{+2}}{P_{+1}^2}
 =g^{(2)}(0)+O(|g_{\mathrm Q}|^2)}
 \label{eq:qpinem-g2-electron-ratio-dp65}
\end{equation}
即可在弱耦合极限中区分亚 Poisson、Poisson 和 bunching 光子统计。这个算例说明 QPINEM 的真正新增信息不是“又出现一套边带”，而是高阶电子边带把光场的下降阶乘矩逐阶映射到电子能谱。
\end{exbox}
